\documentclass[conference]{IEEEtran}

\usepackage{amsmath}
\usepackage{booktabs}
\usepackage{microtype}
\usepackage{xcolor}
\usepackage{url}

\newcommand{\system}{\textsc{Crane-Explain}}
\newcommand{\pending}[1]{\textcolor{red}{[PENDING: #1]}}

\title{Beyond Terminal Events: Evidence-Checked Physical Diagnosis\\
for Natural-Language Robot Navigation Explanations}

\author{Anonymous submission}

\begin{document}
\maketitle

\begin{abstract}
Robot navigation systems expose action results, behavior-tree transitions, maps, commands, and
motion estimates, yet a fluent explanation can still stop at ``recovery was exhausted'' or invent
a physical cause that those signals do not establish. We present a diagnosis-to-language pipeline
that first computes a bounded physical or execution mechanism, links it to exact runtime and source
evidence, records conflicting evidence and unresolved alternatives, and only then realizes and
checks natural language. We distinguish diagnostic correctness from language faithfulness and
measure both against repository-aware agents with matched evidence and tool access. A separately
preserved legacy study evaluates runtime-to-source provenance and factual restraint but is not
reinterpreted as physical diagnosis. The prospective study tests geometric restriction and
command-to-motion mechanisms in simulated land navigation and a focused surface-vessel setting,
including cases with deliberately missing decisive evidence. \pending{Insert independently
annotated supported-diagnostic-success effects, clustered confidence intervals, material-error
risk, and fallback frequency after the protocol is frozen and held-out evaluation completes.}
\end{abstract}

\section{Introduction}

The last message emitted by a navigation stack is rarely its most useful explanation. An action
may abort after recovery exhaustion, but a user or maintainer needs to know what prevented a
successful path or execution. Conversely, telemetry preceding an abort does not license an
unobserved obstacle identity, a unique actuator defect, or a causal claim about wind or waves.
Natural-language generation makes this tension acute: the most readable answer can be the least
auditable one.

We study explanations whose claims are constrained by retained robot evidence. Our central design
separates five stages:
\begin{equation*}
\text{observations and execution evidence}
\rightarrow \text{diagnosis}
\rightarrow \text{checked plan}
\rightarrow \text{language}
\rightarrow \text{final check}.
\end{equation*}
The diagnosis is not free-form reasoning. It is a versioned computation over quantities with
units, frames, time intervals, evidence identifiers, assumptions, and applicability limits. Its
output identifies candidate mechanisms, supporting and contradictory observations, and
alternatives that remain observationally equivalent. The language layer must preserve this
content; if it cannot be verified after one bounded repair, the system falls back to a checked
deterministic rendering.

The paper asks three prospective questions. \emph{Q1:} do physical diagnostic computations plus
runtime/source provenance improve correct and informative failure explanations over a strong
repository-aware agent with equivalent evidence and tools? \emph{Q2:} does final language preserve
the diagnosis, decisive measurements, mechanism, and limitations? \emph{Q3:} does the method give
a mechanism when evidence supports one and narrow its answer when decisive evidence is absent?

This work makes three bounded contributions:
\begin{itemize}
  \item a diagnosis record and checked language path that preserve measurement semantics,
  provenance, conflicts, unresolved alternatives, and permitted causal language;
  \item two navigation-focused diagnostics---geometric restriction and command-to-motion
  discrepancy---with explicit conditions under which they must abstain or qualify; and
  \item a paired evaluation design that separately scores diagnostic correctness, language
  faithfulness, material error, useful coverage, and unnecessary abstention on diagnosable and
  ambiguous cases.
\end{itemize}
We do not claim universal root-cause discovery, human trust, hardware reliability, or that
structured records alone are novel. The empirical scope is simulation-based navigation.

\section{Related Work}

Liu and Brand\~ao generate environment changes that restore feasibility to otherwise unsolved
motion-planning problems~\cite{liu2024environment}. Their work motivates geometric feasibility
checks, but a feasible reconstructed model does not establish successful physical execution, and
local free space does not prove global route absence. Diehl and Ramirez-Amaro learn an explicit
causal model and search for contrastive assignments associated with success~\cite{diehl2022fail}.
Their assumptions clarify why event order alone is not causal evidence.

REFLECT organizes multisensory histories into sensory, event, and subgoal levels to explain robot
failures~\cite{liu2023reflect}. Its reported limitations on low-level control motivate retaining
commands, actuation evidence, and measured motion rather than only scene summaries. HEXAR routes
questions to specialized component explainers and evaluates root-cause identification separately
from incorrect facts~\cite{love2026hexar}. We similarly separate useful diagnosis from unsupported
content, but additionally test whether diagnostic content survives language realization.

Marine motion illustrates the need for semantic restraint: propulsion, current-relative
hydrodynamics, wind, and wave loads are distinct terms in standard craft models~\cite{fossen2021handbook}.
A lateral residual therefore does not uniquely establish ``wave drift.'' More generally, a
navigation feedback field may not mean what its label suggests. In the audited Nav2 Jazzy source,
\texttt{FollowPath.feedback.speed} is computed from the commanded velocity rather than independent
motion~\cite{nav2controller}. Our diagnostic contracts retain command and measured motion as
different evidence types.

\section{Evidence and Diagnostic Method}

\subsection{Evidence boundary}

Robot-visible inputs and evaluator-only truth are physically separated. Robot-visible evidence may
include goals, paths, action/behavior-tree traces, exact source and configuration artifacts,
delivered scans or costmaps, commands, accepted actuation where available, and independently
measured motion. Evaluator truth contains injected conditions, canonical simulator geometry,
hidden force decomposition, expected propositions, and answerability labels. A diagnostic sensor
need not have been consumed by Nav2 to establish a physical condition. Claiming that Nav2 chose an
action \emph{because} it consumed that observation requires a recorded runtime dependency.

We retain four evidence levels: recorded sequence; reconstructed software mechanism; inference
under an explicit model; and controlled-intervention result. Temporal order alone supports only
the first level. Every claim names its evidence IDs and level.

\subsection{Diagnostic result}

A diagnostic result stores the episode and interval; input hashes; measurements with units,
frames, and timestamps; computation name and version; parameters and thresholds; assumptions and
applicability; supporting and contradictory evidence; distinguishable and unresolved alternatives;
the permitted causal-language level; and one fixed-menu next check. Evidence sufficiency is labeled
independently of whether our implementation finds the desired mechanism.

\subsection{Geometric restriction}

The geometric diagnostic compares observed free space with the actual footprint at orientation,
configured safety envelope, and requested path. It reports physical and configured-envelope
clearance separately. A disconnected audited grid can establish no traversable connection only
within that grid, resolution, threshold, and footprint model. Planner failure without such a check
is not proof of infeasibility. When interventions are used, the system executes a new run from a
controlled initial condition; recorded-state playback is not called a counterfactual execution.

\subsection{Command-to-motion discrepancy}

The execution diagnostic aligns desired commands, accepted/applied commands where available,
actuator feedback, and independent pose or velocity. Healthy response bounds are calibrated on
development trials and include delay and inertia. A discrepancy supports an execution-response
mechanism but does not uniquely identify slip, collision, motor failure, or environmental loading.

One concrete terminal-margin check computes
\begin{align*}
m_{\mathrm{return}} &= \tau_{\mathrm{task}} - e_{\mathrm{return}},\\
\Delta e_{\mathrm{settle}} &= e_{\mathrm{settled}} - e_{\mathrm{return}},
\end{align*}
where $\tau_{\mathrm{task}}$ is the declared task tolerance and $e$ is independently measured goal
error. If the action returns under its configured pose and stopped-speed thresholds, but
$\Delta e_{\mathrm{settle}} > m_{\mathrm{return}}$ and the platform settles outside the task
tolerance, the evidence supports an inadequate terminal stopping margin. It does not establish why
the residual motion occurred.

\subsection{Checked answer and final verification}

Answers contain four short sections: \emph{Diagnosis}; \emph{Decisive evidence}; \emph{Failure
chain}; and \emph{Limits and next check}. Measurements are included only when they support the
mechanism. The diagnostic plan is checked before generation, but the final text is also checked.
Unverified clauses are never silently accepted. We retain candidate text, verification outcome,
repair/fallback decision, final text, latency, tokens, and cost where available.

\section{Evaluation Design}

\subsection{Separate legacy evidence set}

Before the redirect, we froze a paired land/Nav2 provenance study over recovery-success and
terminal-abort episodes. Its questions, prompts, models, verifier, rubric, split, and inclusion
rules remain unchanged. Collection was prospectively closed at 33 included episode clusters (17
recovery-success and 16 terminal-abort), below the frozen minimum of 40 and target of 50. The
cohort contains 198 one-shot responses across three conditions. No sealed label or effect estimate
was inspected when collection was closed. The study will be analyzed under its original rubric and
reported as under-target; it is not pooled with or rescored as the diagnostic study.

The legacy study evaluates software-mechanism provenance, factual restraint, and evidence
insufficiency. Its checked condition used deterministic fallback for 46 of 66 responses (69.7\%),
including every recovery-mechanism response. We therefore report candidate and final-output
results separately and describe the final method as predominantly deterministic if this pattern
remains after annotation.

\subsection{Prospective diagnostic comparison}

The new study compares four compact conditions. R is a strong repository-aware agent with the same
robot-visible evidence, exact source/configuration, and the same diagnostic tools and descriptions
available to the proposal. P receives validated diagnostic results plus provenance, checked
planning, generation, and final verification. T deterministically renders the same diagnostic
result. N is a new-arm provenance/checking ablation without physical diagnostic computations; it
does not modify the archived legacy method.

The primary endpoint is \emph{supported diagnostic success} on independently labeled diagnosable
questions: the final answer identifies the required mechanism and contains no material unsupported
or contradicted assertion. Secondary outcomes are diagnostic correctness, language faithfulness,
material-error risk, causal overclaiming, unnecessary abstention, appropriate qualification,
mechanism/location/time identification, evidence coverage and citation correctness, selected next
check, latency/cost, and fallback frequency.

Scenario families contain nominal success, successful compensation, similar terminal symptoms
from different mechanisms, the same mechanism with different outcomes, irrelevant obstacles,
sufficient versus missing evidence, false premises, and out-of-model causes. Shared layouts,
interventions, evidence masks, and question variants stay in one split and one statistical cluster.
Valid unexpected outcomes are retained; induction success and recording validity are reported
separately. Development pilots determine thresholds and inform a new clustered paired power
simulation before questions, labels, prompts, models, stopping rules, and held-out splits freeze.

Two blinded annotators score randomized packets independently and disagreements are adjudicated
before condition keys are joined. Independent reference computations, not the runtime verifier,
provide diagnostic gold labels. Scenario-clustered paired confidence intervals are primary;
response-level paired tests are secondary checks.

\section{Development Case: Terminal Margin}

We first exercised the complete path on a retained RoboBoat development run; it is not a held-out
result. Nav2 returned success at 0.3738~m goal error while independent odometry measured
0.04861~m/s. The exact configuration allowed 0.400~m XY error and 0.050~m/s stopped speed. Against
the declared 0.400~m task tolerance, only 0.0262~m margin remained. The vessel then moved 0.1876~m
and ended at 0.5614~m error. The checked answer therefore leads with an inadequate terminal margin,
connects the configured success to post-result motion, and withholds a unique claim about motor,
hydrodynamic, wind, current, wave, or collision causation.

A paired evidence mask removes only measured speed at return. The diagnostic then reports that the
mechanism is not established and requests synchronized return-speed evidence. This mask verifies
selective behavior in development; it is neither an independent episode nor a confirmatory result.
Matched post-change runs exist as retrospective development context but are excluded from the
ordinary robot-visible packet.

\section{Results}

\begin{table}[t]
\caption{Prospective diagnostic evaluation. Values remain blinded/unavailable until the protocol
is frozen and held-out annotation is complete.}
\label{tab:diagnostic-results}
\centering
\begin{tabular}{lccc}
\toprule
Method & Supported success & Material error & Coverage \\
\midrule
R & \pending{n/N} & \pending{n/N} & \pending{value} \\
P & \pending{n/N} & \pending{n/N} & \pending{value} \\
T & \pending{n/N} & \pending{n/N} & \pending{value} \\
N & \pending{n/N} & \pending{n/N} & \pending{value} \\
\bottomrule
\end{tabular}
\end{table}

\pending{Report legacy annotation/adjudication and its clustered paired estimates separately. Do
not call the early-closed cohort powered, complete under the original stopping rule, or evidence of
physical diagnosis.}

\pending{Report diagnostic pilot-informed sample planning, final cluster counts, exclusions,
induction success, diagnostic-output accuracy, final-language accuracy, risk--coverage curves,
fallback frequency, paired effects with confidence intervals, and all null/negative findings.}

\section{Limitations and Validity}

The diagnostic families are deliberately narrow. A correct result applies only to its retained
model, signals, interval, and assumptions. Delivered observations do not prove controller
consumption; local geometry does not prove global infeasibility; a command--motion residual does
not uniquely identify a physical source. Evaluator truth is unavailable to ordinary methods.

Experiments use simulation. They can test controlled mechanisms and simulation-based ecological
relevance, not hardware reliability or real-world transfer. The study measures factual and
diagnostic outcomes, not human trust. The legacy cohort stopped below its frozen minimum and must
be interpreted through effect sizes and uncertainty rather than significance alone. The
development terminal-margin case was selected retrospectively; only the separately frozen
prospective study can support comparative claims.

Finally, deterministic rendering may outperform or dominate language-model realization. We treat
that as an empirical result, report fallback frequency, and do not attribute templated final
answers to language-model reasoning.

\section{Conclusion}

Useful robot explanations require more than narrating the event that ended an action. Our method
places a validated, bounded diagnosis between robot evidence and language, then checks whether the
final response preserved both the mechanism and its limits. The ongoing evaluation asks whether
that added diagnostic computation improves supported diagnostic success under evidence- and
tool-matched baselines. Regardless of the eventual effect, the artifact preserves the distinction
between what the robot recorded, what software semantics reconstruct, what an explicit model
supports, and what only a controlled intervention can establish.

\bibliographystyle{IEEEtran}
\bibliography{references}

\end{document}
